% =====================================================================
% THIS FILE IS THE SOURCE. Edit it, then rebuild the PDF with:
%     bash tools/compile.sh Unit_02_A_Map_of_Models
% (run from the Course_Rewrite folder)
%
% Unit 2 in the model-first ordering. Introduces the model families,
% splits them into probability models and algorithmic ones, and works
% through three questions: what you want from a model, what you do with
% it, and how you pick. Deliberately shallow on mechanics. Software does
% the arithmetic; this unit is about knowing what to ask for.
% =====================================================================
\documentclass[aspectratio=169]{beamer}
\usetheme{default}
\usefonttheme{professionalfonts}
\usepackage{amsmath, amssymb, amsfonts}
\usepackage{graphicx}
\usepackage{booktabs}
\usepackage{array}
\newcolumntype{L}[1]{>{\raggedright\arraybackslash}p{#1}}
\usepackage{bm}
\usepackage{hyperref}

\mode<presentation> {
    \usetheme{Frankfurt}
    \usecolortheme{default}
}

% == notation shortcuts ==
\newcommand{\xbar}{\bar{x}}
\newcommand{\yhat}{\hat{y}}

% Recurring motifs (color-coded).
\newenvironment{principle}[1]{\begin{block}{\textbf{Principle:} #1}}{\end{block}}
\newenvironment{trap}[1]{\begin{alertblock}{\textbf{Trap:} #1}}{\end{alertblock}}
\newenvironment{practice}[1]{\begin{exampleblock}{\textbf{In practice:} #1}}{\end{exampleblock}}
\newenvironment{interview}[1]{\begin{exampleblock}{\textbf{You may be asked:} #1}}{\end{exampleblock}}
\newenvironment{yourturn}[1]{%
  \setbeamercolor{block title}{fg=white,bg=orange!80!black}%
  \setbeamercolor{block body}{bg=orange!12}%
  \begin{block}{\textbf{Your turn:} #1}}{\end{block}}
\newenvironment{livedemo}[1]{%
  \setbeamercolor{block title}{fg=white,bg=teal!80!black}%
  \setbeamercolor{block body}{bg=teal!12}%
  \begin{block}{\textbf{Live demo:} #1}}{\end{block}}
\newenvironment{llmcheck}[1]{%
  \setbeamercolor{block title}{fg=white,bg=violet!75!black}%
  \setbeamercolor{block body}{bg=violet!10}%
  \begin{block}{\textbf{LLM check:} #1}}{\end{block}}

\title{A Map of Models}
\subtitle{Unit 2: What Is on the Menu, and How to Choose}
\author{Stat 220 \textbar{} Statistics for Data Science}
\date{}

\begin{document}

\begin{frame}
  \titlepage
\end{frame}

% ======================================================================
\section{Why}
% ======================================================================
\begin{frame}{Where we left off}
\begin{itemize}
  \item Unit 1 compared two groups. One estimate, one standard error, one question.
  \item Almost nothing you will be handed at work looks like that. You get a table with forty
        columns and a vague question attached.
  \item This unit is the map: what is available, what each one lets you do, and how to pick
        without guessing.
\end{itemize}
\begin{principle}{the honest starting point}
``Which model should I use?'' has no answer until you say what you want out of it. Most bad
modeling starts with someone skipping that sentence.
\end{principle}
\end{frame}

\begin{frame}{What this unit gives you}
\begin{columns}[T]
\begin{column}{0.5\textwidth}
\textbf{Things to know}
\begin{itemize}\small
  \item the main model families, and what each one assumes
  \item the difference between a probability model and an algorithmic one
  \item the five jobs people hand a model, and how they conflict
  \item which operations a model supports, and which it does not
  \item how sample size and the goal narrow the choice
\end{itemize}
\end{column}
\begin{column}{0.5\textwidth}
\textbf{Mistakes to stop making}
\begin{itemize}\small
  \item picking the algorithm before naming the goal
  \item asking a random forest for a $p$-value
  \item comparing an AIC against a cross-validation score
  \item reading a penalized coefficient as an effect
  \item reaching for flexibility on a few hundred rows
\end{itemize}
\end{column}
\end{columns}
\end{frame}

% ======================================================================
\section{The menu}
% ======================================================================
\begin{frame}{Two kinds of thing get called a model}
\begin{columns}[T]
\begin{column}{0.48\textwidth}
\textbf{A probability model}\\[2pt]
Says where the data came from. It writes down a distribution for the outcome:
\[ y \sim \text{Normal}(b_0 + b_1 x,\; \sigma^2) \]
Everything else follows from that sentence.
\end{column}
\begin{column}{0.48\textwidth}
\textbf{An algorithmic model}\\[2pt]
Says nothing about where the data came from. It gives you a rule that turns inputs into a
prediction, tuned until the predictions are good.\\[4pt]
A random forest is an average of 500 trees. That is the whole claim.
\end{column}
\end{columns}
\vspace{0.4em}
\begin{principle}{this split decides what you are allowed to ask}
A distribution gives you a likelihood, and a likelihood gives you standard errors, $p$-values,
confidence intervals, and AIC. With no distribution there is no likelihood, and none of those
exist. You get uncertainty by resampling or by holding data out instead.
\end{principle}
\end{frame}

\begin{frame}{The probability models}
\begin{center}\footnotesize
\begin{tabular}{L{0.26\textwidth} L{0.26\textwidth} L{0.38\textwidth}}
\toprule
\textbf{Model} & \textbf{Outcome it is for} & \textbf{What it says} \\
\midrule
Linear regression & a number & the mean moves in a straight line, with normal scatter around it \\
Logistic regression & yes or no & the log-odds move in a straight line \\
Poisson regression & a count, or a rate & counts arrive at a rate that depends on the inputs \\
Multinomial & one of several categories & the same idea with more than two outcomes \\
Mixed models & grouped or repeated data & rows inside a group are related, and the model says how \\
Bayesian versions & any of the above & the same models, plus a prior, giving a distribution over the answer \\
\bottomrule
\end{tabular}
\end{center}
\centering\small
All of these are the same move: name a distribution, let the inputs shift its center.
\end{frame}

\begin{frame}{The algorithmic models}
\begin{center}\footnotesize
\begin{tabular}{L{0.26\textwidth} L{0.30\textwidth} L{0.34\textwidth}}
\toprule
\textbf{Model} & \textbf{How it works} & \textbf{What it is good at} \\
\midrule
Decision tree & splits the data on one variable at a time & a rule a human can follow, and read out loud \\
Random forest & averages hundreds of trees grown on resampled data & solid predictions with little tuning \\
Gradient boosting & each tree fixes the last one's mistakes & usually the best raw accuracy on tabular data \\
$k$-nearest neighbours & predicts from the most similar past cases & simple, and surprisingly hard to beat locally \\
Neural networks & layers of weighted sums, fit by gradient descent & images, text, and very large datasets \\
\bottomrule
\end{tabular}
\end{center}
\centering\small
None of these assume a distribution, so none of them hand you a $p$-value.
\end{frame}

\begin{frame}{The awkward middle: penalized regression}
\begin{itemize}
  \item \textbf{Lasso} and \textbf{ridge} look like linear regression with a penalty added for
        large coefficients.
  \item The form is a probability model, so people report the coefficients as if the usual
        machinery applied.
  \item It does not. The penalty deliberately biases every coefficient toward zero, so the
        printed standard errors and $p$-values are not describing what they claim to.
\end{itemize}
\begin{trap}{a coefficient that has been shrunk is not an effect}
Use these to \emph{screen} many candidates down to a few, then refit an ordinary regression on
the survivors if you need to defend a number. Report the screening step when you do.
\end{trap}
\end{frame}

% ======================================================================
\section{What you want from it}
% ======================================================================
\begin{frame}{Five jobs people hand a model}
\begin{enumerate}
  \item \textbf{Predict.} Give me a number for this new case.
  \item \textbf{Explain.} Does $X$ actually move $Y$, and by how much?
  \item \textbf{Screen.} Out of these 200 columns, which few are worth tracking?
  \item \textbf{Decide.} Given what it costs to be wrong each way, what should we do?
  \item \textbf{Deploy.} Run this every night for two years, and be able to explain any single
        output to someone who is unhappy about it.
\end{enumerate}
\begin{principle}{say which one out loud first}
These are five different jobs. A model that is excellent at the first can be useless at the
second, and the choice that makes the fifth easy usually costs you accuracy on the first.
\end{principle}
\end{frame}

\begin{frame}{The jobs pull against each other}
\begin{center}\small
\begin{tabular}{L{0.18\textwidth} L{0.36\textwidth} L{0.36\textwidth}}
\toprule
\textbf{Job} & \textbf{What it needs} & \textbf{What it will cost you} \\
\midrule
Predict  & flexibility, and lots of rows & coefficients nobody can read \\
Explain  & a simple form and honest standard errors & accuracy you could have had \\
Screen   & a method that handles many columns at once & valid $p$-values on the survivors \\
Decide   & calibrated probabilities and a cost table & extra work nobody asked for \\
Deploy   & stability, few moving parts, a clear story & the last few points of accuracy \\
\bottomrule
\end{tabular}
\end{center}
\begin{practice}{a strong thing to say}
``Before I pick anything: is this going in a report, or in production? The answer changes what I
would build.''
\end{practice}
\end{frame}

% ======================================================================
\section{What you do with it}
% ======================================================================
\begin{frame}{The verbs}
Every model gets some of these done to it. Which ones are available is not the same for all of
them.
\begin{itemize}\small
  \item \textbf{Estimate parameters.} Find the slopes, the splits, the weights.
  \item \textbf{Tune hyperparameters.} Settings you choose rather than estimate: tree depth, the
        lasso penalty, $k$ in nearest neighbours.
  \item \textbf{Predict.} Push a new case through and get a number back.
  \item \textbf{Test.} Ask whether a coefficient could plausibly be zero.
  \item \textbf{Put an interval on it.} Around a coefficient, or around a prediction.
\end{itemize}
\begin{principle}{parameters are learned, hyperparameters are chosen}
A parameter comes out of fitting. A hyperparameter you set beforehand, and the honest way to set
it is by cross-validation, never by trying values until the answer looks right.
\end{principle}
\end{frame}

\begin{frame}{Which verbs each family supports}
\begin{center}\footnotesize
\begin{tabular}{L{0.33\textwidth} L{0.26\textwidth} L{0.29\textwidth}}
\toprule
 & \textbf{Probability model} & \textbf{Algorithmic model} \\
\midrule
Parameters mean something & yes, a slope is an effect & no, splits are bookkeeping \\
Standard errors and $p$-values & free from the likelihood & none exist \\
Confidence interval for an effect & yes & not really \\
Prediction interval & from the assumed distribution & from held-out errors \\
AIC or BIC & yes & undefined, no likelihood \\
Out-of-sample error & yes & yes, and it is the only option \\
Hyperparameters to tune & few & many, and they matter \\
\bottomrule
\end{tabular}
\end{center}
\centering\small
The right-hand column is shorter. That is the trade you make for the accuracy.
\end{frame}

\begin{frame}{Two intervals that get confused}
\begin{columns}[T]
\begin{column}{0.48\textwidth}
\textbf{Confidence interval}\\[2pt]
Around a \emph{parameter}. Where is the true slope? Shrinks as $n$ grows, toward zero width.
\end{column}
\begin{column}{0.48\textwidth}
\textbf{Prediction interval}\\[2pt]
Around a \emph{new observation}. Where will the next case land? Stays wide no matter how much
data you collect, because individual cases vary.
\end{column}
\end{columns}
\vspace{0.5em}
\begin{trap}{quoting the wrong one}
``Our model predicts \$420k, 95\% CI \$418k to \$422k'' is almost always a confidence interval
being passed off as a prediction. The honest range for one house is far wider, and the customer
is asking about one house.
\end{trap}
\centering\small Unit 4 takes this apart properly.
\end{frame}

% ======================================================================
\section{Choosing}
% ======================================================================
\begin{frame}{Simple or complicated}
\begin{center}\small
\begin{tabular}{L{0.18\textwidth} L{0.35\textwidth} L{0.35\textwidth}}
\toprule
 & \textbf{Reach for simple} & \textbf{Reach for flexible} \\
\midrule
Goal      & explaining, screening, deploying & predicting \\
Rows      & hundreds & tens of thousands and up \\
Columns   & few, chosen for a reason & many, and you do not know which \\
Shape     & you can say what it should look like & you have no idea, and want the fit to find it \\
Audience  & someone who will challenge a number & someone who only sees the output \\
\bottomrule
\end{tabular}
\end{center}
\begin{principle}{flexibility is bought with sample size}
A complicated model is a bet that you have enough data to pin down everything you just let vary.
On 300 rows that bet usually loses.
\end{principle}
\end{frame}

\begin{frame}{Comparing models by a number}
\begin{itemize}
  \item \textbf{AIC and BIC} trade fit against the number of parameters. They need a likelihood,
        so they only apply to probability models, and only when the models are fit to the same
        rows. Lower is better, and the raw value means nothing on its own.
  \item \textbf{Out-of-sample error} fits on part of the data and scores on the rest. It works
        for every model on this list, including the ones with no likelihood.
  \item Cross-validation is that idea done five or ten times over so the score is less dependent
        on which rows you happened to hold out.
\end{itemize}
\begin{trap}{comparing scores that are not comparable}
An AIC and a cross-validated error are different quantities on different scales. Pick one method
and use it for every candidate, on identical data.
\end{trap}
\end{frame}

\begin{frame}{The two-model habit}
\begin{itemize}
  \item Fit the simplest defensible model and a flexible one. Both, every time.
  \item If they agree, report the simple one. It is easier to defend, cheaper to maintain, and
        less likely to break when the data shifts.
  \item If they disagree, you have learned something specific: there is structure the simple
        model is missing. Go find out what it is.
\end{itemize}
\begin{practice}{a strong thing to say}
``The forest beat the regression by four points, so I looked for what it was using, found the
threshold at 30 days, added that to the regression, and it caught up. I shipped the regression.''
\end{practice}
\end{frame}

\begin{frame}{Your turn \#1}
\begin{yourturn}{name the family and the reason}
\begin{enumerate}\small
  \item 1{,}100 loan applications, 12 clean columns, and a regulator who will ask you to justify
        every rejection.
  \item 6 million ad impressions, 400 columns, and nobody will ever read a coefficient.
  \item 60 trial patients, one treatment variable, and the question is whether the drug works.
  \item 25{,}000 sensor readings where you know the relationship bends but not how.
\end{enumerate}
\end{yourturn}
The family is the easy half. The reason is the half that gets graded.
\end{frame}

\begin{frame}{Your turn \#1: answer}
\small
\begin{enumerate}
  \item \textbf{Logistic regression.} Not because it predicts best, but because ``justify every
        rejection'' means somebody will read a coefficient, and only a probability model gives
        you one you can defend.
  \item \textbf{Boosting or a forest.} Six million rows makes flexibility affordable and nobody
        needs the parameters to mean anything.
  \item \textbf{A $t$-test, or a regression with the treatment indicator.} With $n=60$ and one
        variable of interest, a flexible model would spend your data estimating things nobody
        asked about. Randomization is doing the work here, not the model.
  \item \textbf{Both.} A flexible fit to see the shape, then name it and put it in a regression
        you can explain.
\end{enumerate}
\begin{principle}{}
Three of these four were decided before anyone looked at the data.
\end{principle}
\end{frame}

% ======================================================================
\section{LLM check}
% ======================================================================
\begin{frame}{LLM check: mistake or not?}
\begin{llmcheck}{you asked for help choosing a model}
``With 200 patients and 45 predictors I'd recommend a random forest. It handles nonlinearity and
interactions automatically and gives you feature importances, so you can report the top
predictors as the key risk factors and their $p$-values.''
\end{llmcheck}
\vspace{0.2em}
\centering\small Three separate problems. Find them before the next slide.
\end{frame}

\begin{frame}{LLM check: answer}
\begin{itemize}
  \item \textbf{The ratio.} 200 rows against 45 columns is the corner where flexible models fit
        well and predict badly. A penalized regression is the honest choice.
  \item \textbf{``Automatically'' costs something.} Discovering interactions means spending data
        to find them. With 200 rows you cannot afford the search.
  \item \textbf{A forest has no $p$-values.} There is no likelihood, so the quantity being asked
        for does not exist. And an importance score ranks correlated columns close to
        arbitrarily, so it is not an effect either.
\end{itemize}
\begin{principle}{}
The answer is not wrong about what a forest does. It is wrong about whether this dataset can pay
for one, and about what the output means.
\end{principle}
\end{frame}

% ======================================================================
\section{Consolidate}
% ======================================================================
\begin{frame}{Choosing a model, in order}
\small
\begin{enumerate}
  \item \textbf{Name the job.} Predict, explain, screen, decide, or deploy.
  \item \textbf{Name the outcome type.} A number, a yes or no, a count, a category. This is not a
        judgment call and it removes most of the menu.
  \item \textbf{Decide whether you need parameters that mean something.} If yes, you are in the
        probability column, and that is the end of the discussion.
  \item \textbf{Count what you have against what you are estimating.} Rows, columns, and noise.
  \item \textbf{Fit the simple one first.} It is your baseline and often your answer.
  \item \textbf{Fit a flexible one on the same folds}, compare with one method, and take the
        simpler model when they are close.
  \item \textbf{Report on data you touched once}, and say what else you tried.
\end{enumerate}
\end{frame}

\begin{frame}{The traps, all in one place}
\begin{trap}{the ones that bite most often}
\begin{itemize}
  \item Choosing the algorithm before naming the job.
  \item Asking an algorithmic model for a $p$-value or an AIC.
  \item Reading a penalized or shrunk coefficient as an effect.
  \item Comparing an in-sample criterion against an out-of-sample one.
  \item Tuning a hyperparameter until the answer looks right.
  \item Quoting a confidence interval when the question was about one new case.
  \item Reaching for flexibility you do not have the rows to pay for.
\end{itemize}
\end{trap}
\end{frame}

\begin{frame}{What you should be able to do with software}
\small
Nobody is asking you to memorize syntax. You are expected to know which procedure the question
calls for, what to hand it, and what the output means.
\begin{center}\footnotesize
\begin{tabular}{L{0.44\textwidth} L{0.46\textwidth}}
\toprule
\textbf{You should be able to} & \textbf{What you would reach for} \\
\midrule
Fit a linear or logistic regression & \texttt{smf.ols} / \texttt{smf.logit} \\
Read the coefficient table & \texttt{fit.summary()} \\
Fit a forest or a boosted model & \texttt{RandomForestRegressor}, \texttt{HistGradientBoosting...} \\
Tune a hyperparameter honestly & \texttt{GridSearchCV} over folds you fixed first \\
Compare probability models & \texttt{fit.aic}, on identical rows \\
Compare anything at all & \texttt{cross\_val\_score} with one shared \texttt{KFold} \\
Screen many columns & \texttt{LassoCV}, then read what survived \\
\bottomrule
\end{tabular}
\end{center}
\vspace{0.2em}
\centering\small
If you cannot say what the output means, producing it was worth nothing.
\end{frame}

\begin{frame}{Where we go next}
\begin{itemize}
  \item You now have the map. The rest of Part I walks the parts of it that matter most.
  \item \textbf{Unit 3} takes the linear model seriously: reading a slope, what happens when you
        add a variable, and how to tell whether the fit is real.
  \item Later units come back to the right-hand column, to prediction intervals, and to the
        categorical outcomes this map only names.
\end{itemize}
\begin{principle}{the through-line}
The model is a means. The job you were given decides which one, and saying that job out loud is
most of the work.
\end{principle}
\end{frame}

\end{document}
